% ========================================================
% BÀI TẬP TOÁN 5 - KẾT NỐI TRI THỨC VỚI CUỘC SỐNG
% Bài 72: Ôn tập đo lường
% Định dạng: BTTeX Part III - Trắc nghiệm điền số ngắn
% ========================================================

% Câu 1
\begin{ex}
  Đổi $3\text{ km } 50\text{ m}$ ra mét được bao nhiêu mét?
  \shortans{3050}
  \loigiai{
    $3\,000 + 50 = 3\,050\text{ m}$.
  }
\end{ex}

% Câu 2
\begin{ex}
  Đổi $4{,}5\text{ m}$ ra xăng-ti-mét được bao nhiêu xăng-ti-mét?
  \shortans{450}
  \loigiai{
    $4{,}5 \times 100 = 450\text{ cm}$.
  }
\end{ex}

% Câu 3
\begin{ex}
  Đổi $2$ tấn $50\text{ kg}$ ra ki-lô-gam được bao nhiêu ki-lô-gam?
  \shortans{2050}
  \loigiai{
    $2\,000 + 50 = 2\,050\text{ kg}$.
  }
\end{ex}

% Câu 4
\begin{ex}
  Đổi $3{,}5$ tạ ra ki-lô-gam được bao nhiêu ki-lô-gam?
  \shortans{350}
  \loigiai{
    $3{,}5 \times 100 = 350\text{ kg}$.
  }
\end{ex}

% Câu 5
\begin{ex}
  Đổi $5\text{ m}^2\ 25\text{ dm}^2$ ra đề-xi-mét vuông được bao nhiêu đề-xi-mét vuông?
  \shortans{525}
  \loigiai{
    $500 + 25 = 525\text{ dm}^2$.
  }
\end{ex}

% Câu 6
\begin{ex}
  Đổi $2\text{ ha}$ ra mét vuông được bao nhiêu mét vuông?
  \shortans{20000}
  \loigiai{
    $2 \times 10\,000 = 20\,000\text{ m}^2$.
  }
\end{ex}

% Câu 7
\begin{ex}
  Đổi $1{,}5\text{ ha}$ ra mét vuông được bao nhiêu mét vuông?
  \shortans{15000}
  \loigiai{
    $1{,}5 \times 10\,000 = 15\,000\text{ m}^2$.
  }
\end{ex}

% Câu 8
\begin{ex}
  Đổi $4\text{ m}^3$ ra đề-xi-mét khối được bao nhiêu đề-xi-mét khối?
  \shortans{4000}
  \loigiai{
    $4 \times 1\,000 = 4\,000\text{ dm}^3$.
  }
\end{ex}

% Câu 9
\begin{ex}
  Đổi $2\,500\text{ cm}^3$ ra đề-xi-mét khối. Kết quả nhân 10 là bao nhiêu?
  \shortans{25}
  \loigiai{
    $2\,500 : 1\,000 = 2{,}5\text{ dm}^3$. Nhân 10 được 25.
  }
\end{ex}

% Câu 10
\begin{ex}
  Đổi $3$ giờ $15$ phút ra phút được bao nhiêu phút?
  \shortans{195}
  \loigiai{
    $3 \times 60 + 15 = 195$ phút.
  }
\end{ex}

% Câu 11
\begin{ex}
  Đổi $2{,}5$ ngày ra giờ được bao nhiêu giờ?
  \shortans{60}
  \loigiai{
    $2{,}5 \times 24 = 60$ giờ.
  }
\end{ex}

% Câu 12
\begin{ex}
  Đổi $180$ giây ra phút được bao nhiêu phút?
  \shortans{3}
  \loigiai{
    $180 : 60 = 3$ phút.
  }
\end{ex}

% Câu 13
\begin{ex}
  Một xe tải chở được $3{,}5$ tấn gạo. Đã chở được $1\,500\text{ kg}$. Còn phải chở bao nhiêu ki-lô-gam nữa?
  \shortans{2000}
  \loigiai{
    $3\,500 - 1\,500 = 2\,000\text{ kg}$.
  }
\end{ex}

% Câu 14
\begin{ex}
  Một thửa ruộng thu hoạch được $1{,}2$ tấn thóc. Đổi ra tạ được bao nhiêu tạ?
  \shortans{12}
  \loigiai{
    $1{,}2 \times 10 = 12$ tạ.
  }
\end{ex}

% Câu 15
\begin{ex}
  Đổi $45\,000\text{ m}^2$ ra héc-ta. Kết quả nhân với 10 là bao nhiêu?
  \shortans{45}
  \loigiai{
    $45\,000 : 10\,000 = 4{,}5\text{ ha}$. Nhân 10 được 45.
  }
\end{ex}

% Câu 16
\begin{ex}
  Đổi $5\text{ dm}^3$ ra lít được bao nhiêu lít nước?
  \shortans{5}
  \loigiai{
    $1\text{ dm}^3 = 1$ lít nên $5\text{ dm}^3 = 5$ lít.
  }
\end{ex}

% Câu 17
\begin{ex}
  Một bể chứa được $2\text{ m}^3$ nước. Bể đó chứa được bao nhiêu lít nước?
  \shortans{2000}
  \loigiai{
    $2\text{ m}^3 = 2\,000\text{ dm}^3 = 2\,000$ lít.
  }
\end{ex}

% Câu 18
\begin{ex}
  Đổi $750\text{ mm}$ ra xăng-ti-mét được bao nhiêu xăng-ti-mét?
  \shortans{75}
  \loigiai{
    $750 : 10 = 75\text{ cm}$.
  }
\end{ex}

% Câu 19
\begin{ex}
  Đổi $0{,}8\text{ km}$ ra mét được bao nhiêu mét?
  \shortans{800}
  \loigiai{
    $0{,}8 \times 1\,000 = 800\text{ m}$.
  }
\end{ex}

% Câu 20
\begin{ex}
  Đổi $6\text{ tạ } 5\text{ kg}$ ra ki-lô-gam được bao nhiêu ki-lô-gam?
  \shortans{605}
  \loigiai{
    $600 + 5 = 605\text{ kg}$.
  }
\end{ex}

% Câu 21
\begin{ex}
  Một bao xi măng nặng $50\text{ kg}$. Muốn có $1$ tấn xi măng cần bao nhiêu bao như thế?
  \shortans{20}
  \loigiai{
    $1\,000 : 50 = 20$ bao.
  }
\end{ex}

% Câu 22
\begin{ex}
  Đổi $3\text{ thế kỉ}$ ra năm được bao nhiêu năm?
  \shortans{300}
  \loigiai{
    $3 \times 100 = 300$ năm.
  }
\end{ex}

% Câu 23
\begin{ex}
  Đổi $\frac{3}{5}\text{ m}$ ra xăng-ti-mét được bao nhiêu xăng-ti-mét?
  \shortans{60}
  \loigiai{
    $\frac{3}{5} \times 100 = 60\text{ cm}$.
  }
\end{ex}

% Câu 24
\begin{ex}
  Đổi $\frac{1}{4}\text{ giờ}$ ra phút được bao nhiêu phút?
  \shortans{15}
  \loigiai{
    $\frac{1}{4} \times 60 = 15$ phút.
  }
\end{ex}

% Câu 25
\begin{ex}
  Đổi $\frac{1}{2}\text{ tấn}$ ra ki-lô-gam được bao nhiêu ki-lô-gam?
  \shortans{500}
  \loigiai{
    $\frac{1}{2} \times 1\,000 = 500\text{ kg}$.
  }
\end{ex}
